\subsection{common}\label{common}
Das \verb|common| Modul bildet die gemeinsame Grundlage aller übrigen
Module. Sein Header \verb|common.h| wird praktisch überall eingebunden und
stellt die grundlegenden Typen und Makros sowie eine Reihe von
Hilfsfunktionen (Hex\-Ausgaben, Prüfsummen, Zustandsnamen, Label\-Helfer)
bereit. Über das Makro \verb|UNIT_TEST| wird zwischen dem MCU\-Build (bindet
\verb|hal_port.h| ein) und dem Host\-Test\-Build umgeschaltet; im Testbau ist
\verb|STATIC| leer, sodass dateilokale Funktionen testbar werden.

\subsubsection*{Grundtypen und Makros}
Funktionen melden ihren Status einheitlich über \verb|em_msg|
(Listing \ref{em_msg}). Daneben definiert \verb|common.h| u.\,a.\ die Makros
\verb|MIN| und \verb|MAX|, \verb|ELCNT| (Elementanzahl eines Arrays),
\verb|NL| / \verb|NEWLINE| (Zeilenende) sowie die Union \verb|cPtrAway_u|
zum kontrollierten Entfernen eines \verb|const|.

\begin{listing}
\begin{minted}{C}
typedef enum {
    EM_ERR  = -1,  // Fehler
    EM_OK   =  0,  // Erfolg
    EM_TRUE,       // wahr
} em_msg;
\end{minted}
\caption{Einheitlicher Rückgabetyp \texttt{em\_msg}}
\label{em_msg}
\end{listing}

\subsubsection*{Systemzustände}
Der Ablauf der Funk\-Synchronisation wird über \verb|system_state_e|
(Listing \ref{system_state_e}) beschrieben. Die Abbildung \verb|synca2str|
liefert zu jedem Zustand seinen Namen als Zeichenkette (siehe
\verb|idxa2str()| weiter unten).

\begin{listing}
\begin{minted}{C}
typedef enum {
    SYNC_RESET, BOOT_UP, SLOT, CHANNEL, FREQBAND,
    FREQUENCY_OFFSET, SYNCHRONIZE, SYNCHRONIZE_DOING,
    SYNCHRONIZE_READY, SYNCHRONIZE_ERROR, SYNCHRONIZE_OK,
    SYNC_CNT
} system_state_e;
\end{minted}
\caption{Synchronisations\-Zustände \texttt{system\_state\_e}}
\label{system_state_e}
\end{listing}

\subsubsection*{Index\-zu\-String\-Abbildung}
Für die Umwandlung numerischer Indizes (z.\,B.\ Zustände) in lesbare Namen
dienen die beiden Strukturen aus Listing \ref{idx2str_t}: \verb|idx2str_t|
bildet einen einzelnen Index auf eine Zeichenkette ab, \verb|idxa2str_t|
fasst eine solche Tabelle samt Länge zusammen.
{\sloppy\emergencystretch=3em
\begin{description}
\item[\texttt{idx2str(map, cnt, idx)}] Sucht \verb|idx| in einer Tabelle aus
  \verb|cnt| Einträgen und liefert die zugehörige Zeichenkette bzw.
  \verb|"NA "|, falls nicht gefunden.
\item[\texttt{idxa2str(map, idx)}] Bequemlichkeit für \verb|idxa2str_t|: ruft
  \verb|idx2str()| mit den in \verb|map| hinterlegten Daten auf.
\end{description}
\par}

\begin{listing}
\begin{minted}{C}
typedef struct idx2str_s {
    char    *str;
    uint8_t  idx;
} idx2str_t;

typedef struct idxa2str_s {
    uint8_t    cnt;
    idx2str_t *entry;
} idxa2str_t;
\end{minted}
\caption{Lookup\-Tabellen \texttt{idx2str\_t} / \texttt{idxa2str\_t}}
\label{idx2str_t}
\end{listing}

\subsubsection*{Label\-Helfer}
Die Union \verb|clabel_u| (Listing \ref{clabel_u}) speichert bis zu
\verb|CMD_LEN|$-1$ Zeichen als kurzes Label bzw.\ Kommando und wird u.\,a.\
vom \verb|buffer| (Abschnitt \ref{buffer}) genutzt.
{\sloppy\emergencystretch=3em
\begin{description}
\item[\texttt{clable2type(lbl)}] Bestimmt, ob das Label eine Hexzahl
  (\verb|hexnum|), druckbares ASCII (\verb|ascii|) oder nicht darstellbar
  (\verb|nonasci|) ist (Rückgabetyp \verb|type_e|).
\item[\texttt{clabel2uint(lbl)} / \texttt{str2uint(str)}] Interpretieren das
  Label bzw.\ eine Zeichenkette als vorzeichenlose Ganzzahl.
\end{description}
\par}

\begin{listing}
\begin{minted}{C}
typedef union {
    uint32_t cmd;        // als 32-Bit-Kommando
    char     str[CMD_LEN]; // oder als bis zu CMD_LEN Zeichen
} clabel_u;
\end{minted}
\caption{Label\-/Kommando\-Union \texttt{clabel\_u}}
\label{clabel_u}
\end{listing}

\subsubsection*{Hex\-Ausgabe und Dumps}
{\sloppy\emergencystretch=3em
\begin{description}
\item[\texttt{to\_hex(out, out\_size, buffer, buffer\_size, write\_asci)}]
  Formatiert \verb|buffer| als Hexdump in den Zeichenpuffer \verb|out|: je
  Zeile die Adresse, $16$ Hexbytes (mit Extra\-Abstand nach $8$) und optional
  die ASCII\-Darstellung. Liefert die Anzahl geschriebener Zeichen.
\item[\texttt{print\_buffer(buffer, size, header)}] Gibt einen Hexdump
  ($8$ Bytes je Zeile mit ASCII\-Spalte) direkt über \verb|printf| aus.
\item[\texttt{int2hchar(nr)}] Wandelt einen Wert in das zugehörige
  Hexziffer\-Zeichen (\texttt{0}--\texttt{9}, \texttt{A}--\texttt{F}) um.
\end{description}
\par}

\subsubsection*{Berechnungen}
{\sloppy\emergencystretch=3em
\begin{description}
\item[\texttt{common\_crc16(data\_p, length)}] Berechnet eine CRC\-16 nach
  X.25 (Polynom \verb|0x8408| reversed, Startwert \verb|0xFFFF|, XOR\-out
  \verb|0xFFFF|). Das Ergebnis ist kompatibel zur \verb|x-25|\-Definition aus
  Pythons \verb|crcmod|.
\item[\texttt{csss2uint32(cycle, slot, sSlot)}] Packt Cycle ($16$ Bit), Slot
  und Sub\-Slot in einen \verb|uint32_t| (\verb!cycle | slot<<16 | sSlot<<24!).
\item[\texttt{modulo\_sub(slot, oSlot, modulo)}] Liefert die modulare Differenz
  zweier Slots. Im Quelltext als noch nicht getestet markiert.
\end{description}
\par}

\subsubsection*{Hardware\-Bezug}
{\sloppy\emergencystretch=3em
\begin{description}
\item[\texttt{board\_get\_unique\_id(id, max\_len)}] Liest die 96\-Bit
  ($12$~Byte) eindeutige Chip\-ID des STM32L476 (Adresse \verb|UID_BASE|) und
  füllt einen größeren Puffer mit Nullen auf. Im \verb|UNIT_TEST|\-Build durch
  eine deterministische Ersatzimplementierung ersetzt.
\item[\texttt{in\_interrupt()}] Platzhalter; liefert derzeit stets
  \verb|false|.
\end{description}
\par}
